\documentclass[11pt,a4paper]{article}
\usepackage[T1]{fontenc}
\usepackage[margin=27mm]{geometry}
\usepackage{lmodern}
\usepackage{microtype}
\usepackage{amsmath,booktabs,tabularx}
\usepackage[english]{babel}
\usepackage{csquotes}
\usepackage[style=apa,backend=biber,natbib=true]{biblatex}
\addbibresource{references.bib}
\usepackage{tikz}
\usetikzlibrary{positioning,arrows.meta}
\usepackage[hidelinks]{hyperref}
\title{Supplementary material: AI-Assisted Music Creation and Emotion Regulation}
\author{}
\date{}
\begin{document}
\maketitle
This supplement documents the theoretical rationale, proposed protocols and synthetic design diagnostics accompanying AIM-ER. No participant observations are reported. All recruitment and publication gates remain prospective.

\section*{S1. Neuropsychological rationale}
\subsection{Neuropsychological mechanisms: a multilevel account}
AIM-ER distinguishes four interacting neurocognitive functions: auditory-affective appraisal, reward learning, autobiographical retrieval and executive regulation. These are functional domains rather than a one-to-one mapping between psychological constructs and brain regions. The available neuroscience concerns primarily music perception and established emotion-regulation tasks, not the entire AI-assisted creative cycle.

\paragraph{Auditory-affective appraisal.} Musical structure provides temporally organised information from which listeners infer expressive features and form expectations. A distributed network involving auditory regions and structures implicated in affective processing supports musical emotion \citep{koelsch2014}. In an AI-assisted session, successive candidate outputs allow a person to compare perceived expression with the intended affective quality. A discrepancy may elicit revision, rejection or acceptance of an unexpected interpretation. Measure perceived emotion, experienced emotion and intended emotion separately; agreement between them is not required for adaptive regulation.

\paragraph{Reward and selection.} PET and fMRI evidence links anticipatory and peak pleasurable musical experiences to distinguishable striatal responses \citep{salimpoor2011}. Pharmacological modulation of dopaminergic function also changes musical reward experience \citep{ferreri2019}. These findings support a candidate account of repeated selection: anticipation of a satisfactory candidate motivates another attempt, whereas an accepted output provides consummatory pleasure. However, reward-related repetition does not guarantee goal attainment. A creator may continue generating because the next candidate could be more rewarding even after the initial emotional objective has been reached. Specific musical anhedonia further demonstrates that sensitivity to music-related reward varies independently of general reward responsiveness \citep{mas2014}. Measure musical reward sensitivity as a potential moderator, and distinguish anticipated enjoyment, achieved enjoyment, generation count and regulatory outcome. Neither behavioural engagement nor autonomic arousal is a direct measure of dopamine.

\paragraph{Autobiographical memory and narrative reconstruction.} \citet{janata2009} associated autobiographical salience during familiar music listening with activity in medial prefrontal regions. In AIM-ER, a newly generated piece can be deliberately constructed around a pre-existing memory, rather than retrieve memories solely through prior associations with familiar music. Repeated listening may subsequently establish a new association between the artefact and the experience. The two routes require distinct measures: initial autobiographical relevance, later familiarity, memory vividness and change in appraisal. Memories need not be disclosed to establish these variables. Personal resonance is not necessarily accurate autobiographical reconstruction, and intense recollection may either facilitate reflection or sustain distress.

\paragraph{Executive regulation and self-distancing.} Reappraisal recruits distributed cognitive control and semantic processes that alter emotional responding \citep{ochsner2012}. Self-distancing can support reflection without perseverative immersion \citep{kross2011}. A candidate AI-specific affordance is the rapid transformation of a personal narrative into alternative perspectives, voices or musical settings. A useful experiment compares first-person and observer-perspective renderings of equivalent non-sensitive content, with matched duration and independently rated semantic fidelity. Assess changes in psychological distance and interpretation before assessing delayed distress. An observer perspective is a manipulation, not proof of distancing; increased distance might also reflect disengagement or avoidance.

\paragraph{Evidence transfer.} The auditory and reward pathways have direct evidence for music perception, not generation; use matched playback and prospective reward ratings. Autobiographical retrieval findings concern familiar music; distinguish initial narrative relevance from later learned familiarity. Reappraisal and affect labelling have direct evidence in non-musical tasks; compare equivalent content across perspective and labelling manipulations. The proposed link from these operations to regulation in a generator workflow remains a hypothesis, and a neural correlate alone cannot distinguish the competing psychological explanations.

\paragraph{Psychophysiological measurement.} Electrodermal activity and cardiac measures may index changes in autonomic activation during generation, playback and recovery. Their interpretation depends on movement, respiration, sound level, attention and baseline activation. For example, compare matched silent decision periods and playback intervals, record respiration if heart rate variability is analysed, and prespecify artefact handling and minimum usable recording duration. EEG could examine auditory or attentional responses in a separate controlled study; neither scalp spectral power nor a frontal asymmetry score uniquely identifies emotional regulation, agency or reward. Neuroimaging and pharmacological methods are optional mechanistic extensions, not prerequisites for initial behavioural validation.

\subsection{Strength and limits of transferable evidence}
The musicking meta-analysis discussed above reported a trim-and-fill estimate of .40 \citep{peters2024}; with only eight studies, nonsignificant heterogeneity does not establish common effects across activities. Its pooled estimate is unsuitable for planning an incremental effect of AI-assisted generation. A broader synthesis of 104 randomised trials involving 9,617 participants found favourable effects of music interventions on physiological and psychological stress-related outcomes \citep{dewitte2020}. Stress reduction is not interchangeable with goal attainment, and these pooled results do not isolate active generative creation or the proposed goal-by-affordance interaction. A recent mini-review of responses to generated music distinguishes perceived authorship, affective response and aesthetic judgement, while identifying uncertainty about enduring personal meaning \citep{hesun2026}. This listener-focused evidence motivates separating source expectations from perceived resonance in AIM-ER; it does not demonstrate a regulatory advantage of co-creation. Neuroimaging studies establish plausible biological substrates of specific musical or regulatory operations, whereas evidence from AI co-creation studies concerns the affordances of the interface. Neither line alone establishes the incremental regulatory benefit of generative systems. A recent critical review of algorithmically personalised and AI-generated music found very low certainty across assessed physiological, behavioural and clinical outcome domains, with few studies of actual generative models and no independent replication of specific findings \citep{jin2026}. Distinguish emotional conveyance from induction and regulation, and do not transfer results from passive algorithmic delivery to active generation. The strongest test of AIM-ER therefore combines controlled manipulation of a creative operation with temporally ordered psychological outcomes and an active comparator.

For theory evaluation, distinguish three levels of evidential proximity. Evidence about interaction with generative music systems supports specific claims about the availability and experience of creative operations; it does not by itself establish regulatory benefits. Evidence from music perception and non-musical emotion regulation supports the plausibility of component processes under the conditions studied, but its transfer to AI-assisted sessions remains an empirical question. The goal-specific interaction, the temporal coupling of mechanisms and outcomes, and the proposed adverse trajectories are AIM-ER predictions requiring their own tests. These levels describe directness to the target process, not a formal certainty rating. Methodological quality and replication must be evaluated separately. An observed component effect cannot establish the complete process model, and a favourable outcome alone cannot identify its mechanism.

\paragraph{Neuropsychological predictions and limits.} The articulation pathway is most directly relevant to H2 and H5, the appraisal and executive-regulation pathways to CP and H6, and reward and selection to H4 and H10. H3 concerns attribution and ownership rather than a specific neural substrate. H7--H9 concern social feedback, expertise and strategy switching; their neural implementation is not specified here. The decisive behavioural tests compare assigned operations and outcomes; neural or autonomic measures can provide convergent evidence but cannot substitute for them.


\section*{S2. Competing accounts and evidence}
\subsection{Relation to competing accounts}
The distinction among goals, strategies, tactics and mechanisms predates AIM-ER \citep{baltazar2016,baltazar2019}. The framework therefore does not claim a new general principle of regulatory fit. Table~\ref{tab:competing} contrasts explanatory scope and empirical commitments. The domain extension is the production of a stimulus through iterative human--generator decisions; a proposed mechanistic extension is the interaction of specific semantic operations with instructed goals, beyond preference, pleasure or global agency. This extension is not exclusive to AI and could fail even if existing regulation theories remain useful.

\begin{table}[htbp]
\centering\small
\caption{Explanatory scope and discriminating tests.}
\label{tab:competing}
\begin{tabularx}{\linewidth}{@{}>{\raggedright\arraybackslash}p{.26\linewidth}>{\raggedright\arraybackslash}p{.27\linewidth}>{\raggedright\arraybackslash}X@{}}
\toprule
Account & Unit and mechanism & Incremental prediction and test \\
\midrule
Process regulation \citep{gross2015} & Regulation episode; identification, selection and implementation & Manipulate semantic editability within the same nominal strategy; an interaction conditional on the goal is more specific than a general strategy effect. \\
Musical affect self-regulation \citep{baltazar2019} & Music engagement; strategies and induction mechanisms & Hold exposure fixed with yoked playback; an advantage of active goal-congruent revision would implicate production operations beyond receptive engagement. \\
Regulatory fit \citep{higgins2000} & Goal pursuit; fit of orientation and means & Contrast preservation versus perspective transformation at comparable ratings of liking and perceived fit; a content-specific pattern is more restrictive than general value from fit. \\
Regulatory flexibility \citep{bonanno2013} & Person across contexts; repertoire, monitoring and feedback & Manipulate availability of a targeted operation within a session; subsequent switching is a separate prediction rather than a necessary immediate benefit. \\
Self-determination \citep{ryan2000} & Activity; autonomy and competence & At matched decision budgets compare semantic options; global autonomy alone does not require the predicted crossover. \\
\bottomrule
\end{tabularx}
\end{table}

Fidelity preservation should favour articulation, whereas perspective transformation should favour reinterpretation. Attenuation under yoked playback would support a production pathway; replication with equivalent non-AI editing would indicate a domain-general mechanism. A precise absent or reversed matching advantage would weaken this extension, without refuting existing accounts.

\subsection{Recent direct evidence and overlap}
\citet{he2026} manipulated automation in a music co-creation task and examined task load, ownership and state agency across expertise groups. This precedes the control and expertise propositions. \citet{zhang2026} randomly assigned young adults with high alexithymic traits to two songwriting workflows. The accepted abstract reports advantages of user-directed creation for ownership, perceived expressive authenticity and subsequent unaided articulation, but no detectable between-workflow difference in negative-affect change. The latter is not evidence of equivalence. Assignment varied revision, engagement and time as well as control; the intermediate measures were not randomised, and neither a non-AI nor a non-songwriting control was included. The full article was not yet available at the evidence update on 26 November 2026.

\begin{table}[htbp]
\centering\small
\caption{Overlap of AIM-ER propositions with direct co-creation evidence.}
\begin{tabularx}{\linewidth}{@{}>{\raggedright\arraybackslash}p{.1\linewidth}>{\raggedright\arraybackslash}X>{\raggedright\arraybackslash}X@{}}
\toprule
 & He (2026) & Zhang (2026; abstract) and incremental test \\
\midrule
H3 & Ownership and agency covary; authorship distinct. & Workflow affects ownership; test separate, temporally ordered attribution and ownership measures. \\
H4 & Automation affects agency; conceptual replication. & Workflow bundle differs in control; isolate editable semantic options at equal decision budgets. \\
H5 & No trait clarity prediction tested. & Restricted high-alexithymia group; test continuous trait difficulty and distinct state clarity across its range. \\
H8 & Expertise moderates automation responses; replication. & No expertise interaction established in abstract; test continuous expertise with baseline tool familiarity. \\
\bottomrule
\end{tabularx}
\end{table}

\citet{swarbrick2026} report exploratory single-session AI songwriting with musicians and nonmusicians; the conference abstract suggests expertise-related differences in perceived experience, not a randomised effect on regulation. \citet{zhangliu2026} describe goal setting, monitoring and revision in three music students using generative AI. This qualitative microanalysis precedes the proposed iterative architecture, but concerns self-regulated learning rather than instructed emotional goals. The goal--affordance contrast therefore concerns a narrower affective prediction, not the discovery of goal-led revision or expertise differences.

Steering and expressive communication precede AIM-ER \citep{louie2020,louie2022}; PAGURI examines interaction with text-to-music models \citep{ronchini2025}. \citet{zhangpan2026} analyse 304 prompts from 24 PAGURI users and propose control, ownership and expertise pathways overlapping H3, H4 and H8, without measuring these psychological mediators. Distinguish model coverage from editability and lexical continuity from semantic goal attainment. Independently randomise goals and affordances with matched model, decision budget, duration and exposure. Rival explanations include:

\begin{table}[htbp]
\centering\small
\caption{Competing explanations for a matching contrast.}
\begin{tabularx}{\linewidth}{@{}>{\raggedright\arraybackslash}p{.25\linewidth}>{\raggedright\arraybackslash}X>{\raggedright\arraybackslash}X@{}}
\toprule
Alternative & Expected pattern & Design check \\
\midrule
Acoustic or lyric quality & One workflow preferred across goals & Blind ratings of audio quality and semantic fidelity; match candidate pools.\\
Expectancy and demand & Apparent benefits track prior expectations & Balanced instructions; elicit expectations before assignment.\\
Personalisation and exposure & Effect survives passive playback & Yoked audio and equivalent listening durations.\\
Effort or task satisfaction & Effects track effort and decisions regardless of goal & Equal decision budget; separately measure difficulty, satisfaction and effort.\\
General creativity & Equivalent matching effect without AI & Non-AI interface offering comparable semantic operations.\\
\bottomrule
\end{tabularx}
\end{table}
These variables are diagnostics or secondary endpoints, not post-treatment covariates in the primary intention-to-treat contrast.


\section*{S3. Identification and illustrative models}
\subsection{Identification and computational illustrations}
Consider a randomised binary assignment $Z$ to a creative affordance, a proximal mediator $M$ (for example state agency) and a continuous outcome $Y$ oriented so that larger values are favourable. A linear illustrative structural model is
\begin{align}
M &= aZ+U,\\
Y &= \tau Z+bM+cU+\epsilon,
\end{align}
where $U$ is an unobserved mediator--outcome common cause and both disturbances are independent of assignment. For any real $\delta$, define
\[
 b'=b+\delta,\quad c'=c-\delta,\quad \tau'=\tau-a\delta .
\]
Substitution yields the same $M$ and $Y$ for every realisation of $(Z,U,\epsilon)$. Consequently, even randomisation and perfect observation of the mediator do not distinguish these structural models. The total effect $\tau+ab$ is invariant, while the path-specific effect $ab$ changes. This counterexample uses an additive linear model with no treatment--mediator interaction; it demonstrates non-identification in this class, not that every mediation design fails. Randomising a mediator-relevant intervention can test a component-level effect; attribution to the mediator still requires assumptions about other pathways, intervention versions and mediator--outcome confounding.

\paragraph{Goal--affordance contrast.} Let $\mu_{ga}$ denote mean goal attainment under random assignment to goal instruction $g$ and interface $a$, both coded as zero (articulation or fidelity preservation) and one (reinterpretation or perspective transformation). The prespecified equal-cell-weight matching contrast is
\[
 \psi=\frac{\mu_{00}+\mu_{11}-\mu_{01}-\mu_{10}}{2}.
\]
It equals half the conventional difference-in-differences interaction and is positive when matched conditions have higher average attainment than mismatched conditions. Randomise the two factors independently, retain all assigned participants in the estimand, and assess measurement comparability across goal instructions. A positive interaction does not imply that each matched cell individually exceeds its corresponding mismatched cell; report all four means and both simple effects. Confidence intervals and sample size planning must target $\psi$, not the main effect of the interface. Under independent equal-size cells with common residual variance $\sigma^2$ and $n$ participants per cell, $\operatorname{Var}(\hat\psi)=\sigma^2/n$. This planning approximation excludes attrition, clustering and heterogeneous residual variances.

\paragraph{Immediate and delayed effects.} A minimal illustrative response model is
\[
 L=\rho R-\lambda Q,\qquad 0\leq\rho\leq1,\quad \lambda Q\geq0,
\]
where $R$ denotes immediate relief, $L$ delayed relief, $\rho$ persistence and $\lambda Q$ an incremental delayed cost of perseverative thought. For positive immediate relief, delayed deterioration occurs when $\lambda Q>\rho R$. This algebraic example is a possible response trajectory, not a fitted model; other causes of delayed affective change are omitted. The empirical diary model in the validation programme allows additional predictors and participant-level variation.

\paragraph{Reproducible synthetic examples.} A companion Python script illustrates these three properties using arbitrary standardised parameters and a fixed pseudo-random seed. In the first example, models with indirect effects of 0.30 and 0.00 both imply a total effect of 0.40 and produce identical observations up to floating-point rounding. In the factorial example, 1000 simulated observations per cell with $\psi=0.40$ produce $\hat\psi\approx0.382$. Finally, immediate relief of 0.60 and persistence of 0.50 imply delayed relief of 0.30 without perseverative cost and $-0.50$ when that cost is 0.80. These scenarios illustrate structural implications and estimator behaviour; parameter values are not estimates from the target population and should not be used as evidence for efficacy or as an empirical basis for sample size.


\section*{S4. Measurement architecture}
\subsection{Measurement principles}
Table~\ref{tab:measures} separates person-level characteristics, session operations, proposed mediators and outcomes. Instruments are candidates; selection depends on population, language, permissions, recall period and psychometric performance. Brief bespoke state items require cognitive interviewing and pilot evaluation, and should not be described as validated scales. Trait alexithymia may be assessed with an authorised version of the TAS-20 \citep{bagby1994}; the Healthy-Unhealthy Music Scale \citep{hums2015} is an adjacent measure of music-use patterns, not a validated measure of AI-assisted creativity.

\begin{table}[htbp]
\centering\small
\caption{Candidate measurement architecture for AIM-ER.}
\label{tab:measures}
\begin{tabularx}{\linewidth}{@{}p{0.25\linewidth}p{0.19\linewidth}X@{}}
\toprule
Construct & Timing & Operational definition \\
\midrule
Goal and initial affect & Before session & Goal category, distress, valence and arousal; distinguish desired from felt emotion. \\
Trait characteristics & Baseline & Alexithymia, regulatory flexibility, musical expertise and prior AI experience. \\
Creative control & During session & Assigned interface; available and exercised editing, selection and revision operations. \\
Agency and ownership & After operations & Separate state ratings of perceived influence and felt belonging of the output. \\
Personal resonance & After exposure & Perceived correspondence between artefact and intended personal meaning. \\
Clarity and reappraisal & Before and after & Ability to identify emotions and change in interpretation; distinguish from liking. \\
Rumination and avoidance & After; follow-up & Repetitive negative thought and avoidance of relevant tasks or experiences. \\
Affective outcomes & Before, after, follow-up & Distress, valence and arousal on consistent scales; goal attainment separately. \\
Flexibility and functioning & Follow-up & Switching, stopping, return to intended activities; opportunity to switch recorded. \\
\bottomrule
\end{tabularx}
\end{table}


\section*{S5. Detailed validation protocols}
\subsection{Intensive longitudinal study}
A prospective diary design can separate between-person differences from within-person fluctuations. Collect baseline ratings immediately before self-initiated sessions, process measures at the end and follow-up ratings at a fixed interval (for example 24 hours), together with comparable non-creation observations where feasible. Record other regulatory activities, elapsed time, session duration, generator and interface version, and whether the participant revisits the output. Avoid capturing raw prompts or audio by default.

For H6, let $Q_{it}$ denote state rumination assessed after immediate relief and before the delayed outcome, and $\bar Q_i$ the participant mean in a retrospective explanatory model. A candidate mixed model for delayed relief is
\begin{align}
L_{it}={}&\alpha+u_i+\beta_W(Q_{it}-\bar Q_i)+\beta_B\bar Q_i\nonumber\\
&+\gamma R_{it}+\delta D_{it0}+\boldsymbol{\theta}^{\top}C_{it}+e_{it}.
\end{align}
The within-person coefficient $\beta_W$ is predicted to be negative. The primary marginal regression uses participant-cluster CR2 uncertainty and contrast-specific Satterthwaite degrees of freedom (two-sided $\alpha=.04$); the mixed model is a sensitivity analysis. The candidate .04 threshold was selected in synthetic design development to reduce null rejection under the examined MAR scenario (.037 at this threshold); it is not a universal correction for clustering or missingness. Under the examined MNAR scenario the rejection rate remained .067. Independently recalibrate the prespecified analysis and decision rule over plausible designs before confirmatory use; report the nominal threshold and uncertainty explicitly. Interpret the coefficient as an adjusted association, not a causal effect. Where follow-up timing varies, specify a window, record actual delay and test robustness to deviations. Missing follow-up may depend on distress or disengagement; compare observed and missing cases on available baseline variables and conduct sensitivity analyses for departures from missing at random. Estimate uncertainty with participant-level dependence respected. Decompose session-varying predictors into within- and between-person components; add $u_{Qi}(Q_{it}-\bar Q_i)$ with correlated random intercept and slope, prior-session distress and rumination, and residual correlation $\phi^{|t-s|}$ for equally spaced sessions (continuous-time decay for irregular intervals). A sensitivity model predicting $D_{it2}$ with $D_{it0}$ and $D_{it1}$ has the opposite sign on the rumination coefficient under the equivalent complete-data linear parameterisation. Examine lagged rumination and follow-up timing; for prospective prediction use only past observations to construct person means. Compare MAR likelihood, estimated observation weights and pattern-mixture shifts with participant-bootstrap intervals; under MNAR, treat conclusions sensitive to plausible shifts as inconclusive. Shared baseline measurement error and conditioning on post-session relief limit interpretation; the coefficient is not a causal effect.

\subsection{Candidate protocol for the central prediction}
The primary validation study is a parallel-group experiment with four equally allocated conditions. The target population comprises adults able to provide informed consent, understand the study language and complete a brief digital music task. Recruitment should include people with different levels of musical and generative AI experience, rather than only established creator-community members. The initial study should use a standardised fictional or mild, non-sensitive emotional scenario. It does not recruit on the basis of a clinical diagnosis or require personal disclosure. Record eligibility, refusal and recruitment channel before randomisation. The unit of randomisation and primary analysis is the participant, with one assigned session per participant. Determine the recruitment target only after feasibility work, measurement evaluation and joint-power recalibration for CP; the synthetic candidate sizes below are not approved targets.

After consent and baseline affect, experience and expectations are recorded, use a computer-generated allocation sequence concealed until eligibility and baseline assessment are complete. Independently assign goal (articulation or reinterpretation) and semantic affordance (fidelity preservation or perspective transformation). Provide a brief common orientation to the interface. In the articulation condition, ask participants to express the emotional meaning of the scenario accurately. In the reinterpretation condition, ask them to explore a different plausible interpretation of the same scenario. Both interfaces should offer a comparable number of decisions, with one designed to preserve the original meaning and the other to permit a change of perspective. Neither interface should be described as preferable. Use the same input scenario pool, generator version, candidate-generation budget, response format and session time across conditions. A feasibility study must establish that both interfaces permit meaningful user input and implement their intended semantic operations. Participants cannot be blinded to the task instructions; conceal the predicted matching pattern and blind independent outcome coders to assignment where feasible.

The apparatus is a computer-based research interface with controlled access to an approved generation service and standardised audio playback. Record the model and interface version, generation latency, attempted and successful operations, selection or rejection, technical failures and early termination with pseudonymous identifiers. Use identical playback settings within each session and a safe participant-adjustable volume, recording adjustments where feasible. Freeze the software and task materials before confirmatory recruitment. If the generator changes during a study, document the affected sessions and evaluate the impact according to a prespecified version-change rule rather than silently combining incompatible implementations.

The assessment sequence is baseline, assignment and instructions, pre-task expectancy, creative operations, immediate goal attainment, secondary process and affect ratings, and a delayed follow-up. Obtain the primary attainment rating before detailed questions about potential mechanisms to reduce demand priming. Obtain it at the scheduled end of the session even if no usable artefact was produced, because the endpoint concerns the assigned goal rather than product quality. An absent artefact is a process outcome, not a missing participant; a missing rating is not automatically scored zero. Artefact-specific ratings may be undefined when no artefact exists and require a separately specified estimand or descriptive analysis. Delayed functioning and affect are secondary for CP. The primary analysis retains assigned groups and does not exclude participants on the basis of manipulation-check responses, achieved editing, artefact quality or session completion.

Before confirmatory use, run cognitive interviews and feasibility sessions to evaluate instruction comprehension, semantic separation, rating interpretation, completion and technical stability. Have independent raters assess semantic preservation and perspective transformation using a prospective rubric on approved fictional materials, with agreement and uncertainty reported. Match or constrain acoustic features and candidate budgets during design rather than adjusting for achieved quality after randomisation. Prespecify acceptable manipulation fidelity, assessment completion, software reliability and measurement comparability in the feasibility protocol. Failure to meet a criterion calls for revision and another feasibility assessment, not selective exclusion of participants from a confirmatory analysis. The primary factorial test addresses matching within the generative setting; equivalent non-AI production and yoked playback are separate comparisons, not conditions implicitly included in the four-cell design.

\subsection{Experimental decomposition}
A first experiment can isolate the effect of adding a generated musical rendering to a common expressive task. Randomise adults after a standardised mild, non-sensitive writing task to (A) reading their own text, (B) listening to a musical rendering of that text generated using a standardised procedure, or (C) listening to a duration-matched, non-personal musical excerpt. Use the same time budget, comparable playback levels and matched instructions. Condition B versus A estimates the incremental effect of the rendering package, while B versus C addresses personalisation combined with text-content differences; neither isolates AI from musical composition generally. An additional human-composed rendering arm is necessary to isolate differences associated with the production method, and requires appropriate matching.

For H4 a separate randomised comparison assigns an editable versus selection-only interface while holding model, input material, exposure time and available candidate outputs as constant as feasible. State agency is assessed before any social sharing. Report the intention-to-treat contrast and manipulation checks separately; do not condition the primary estimate on achieved editing behaviour. Measure expectation, quality and perceived effort, but avoid adjusting the total effect for post-randomisation variables.

For a distress-reduction experiment addressing H1, predefine the population and comparator, and randomise at the point of eligibility. The primary estimand is the mean difference in immediate distress under assignment to AI creation versus neutral activity, with baseline distress adjustment for precision. A two-arm analysis can use
\begin{equation}
D_{i1}=\alpha+\tau Z_i+\gamma D_{i0}+\boldsymbol{\theta}^{\top}X_i+e_i,
\end{equation}
where $Z_i=1$ denotes AI creation, and $\tau<0$ favours creation. Prespecify covariates, a smallest effect of interest, a two-sided confidence interval, attrition handling and a power calculation tied to the primary endpoint. For H2 the primary endpoint is resonance; H1 and H2 thus need distinct power calculations or a justified multiple-testing plan. Randomise H1 after baseline; use an active non-AI comparator in a subsequent trial. For H2 add a spoken rendering when feasible to separate modality and performance from composition, use a common resonance scale, and obtain blinded fidelity and quality ratings. Record expectations before assignment and after the instruction separately. Use a prespecified multiplicity correction for additional confirmatory contrasts, with exploratory mechanism analyses labelled accordingly. Tests of mediation need temporal separation and explicit assumptions about mediator--outcome confounding; where possible manipulate the proposed mechanism directly.

\subsection{Additional proposition-specific contrasts}
For H3 measure attributed contribution before ownership, with non-overlapping items; compare the adjusted association with the reverse temporal order, controlling prior experience and effort without causal attribution. H5 uses a trait difficulty-identifying-emotions measure and distinct state clarity and unaided articulation outcomes; account for baseline clarity, negative affect and depressive symptoms. H7 assigns opt-in participants to standardised supportive or neutral responses of comparable length and source after a fictional task, assesses validation immediately and goal attainment at a fixed follow-up, and does not imply that actual peers wrote the messages. H8 uses continuous musical expertise and its interaction with randomised interface assignment; assess digital audio workstation familiarity, AI experience and musical identity at baseline. Report conditional contrasts across the observed range without dichotomising expertise. H9 defines an unsuccessful session by a prespecified attainment threshold at the decision point, records available alternatives and separates repeating, switching and stopping (including time-limited and intentional stopping). Analyse competing transitions among eligible decisions, check threshold sensitivity and interpret the conditional association as predictive rather than causal. H10 randomises fixed-version tools with comparable tasks. Offer the final standardised session independently of intermediate attendance. The primary contrast is the between-arm difference of mean first-to-last enjoyment changes. Regress change on first-session enjoyment separately by arm and standardise to each arm's own first-rating distribution, with influence-function uncertainty. This post-assignment rating is not a pretreatment covariate. Under MAR and correct outcome regression, retain the assigned population; model missing first ratings separately and assess arm-specific shifts for missing endpoints. Repeated-session trends, novelty and regulation are secondary outcomes; novelty mediation is not identified.

\section*{S6. Synthetic calibration and sample-size sensitivity}
\subsection{Precision, power and staged validation}
Let $d$ denote the two-arm mean difference in absolute value, divided by the marginal outcome standard deviation, and $d_\psi=\psi/\sigma_Y$ the corresponding standardised matching contrast (half the conventional difference-in-differences). Under equal allocation, independent observations and a common residual variance $\sigma_Y^2(1-R^2)$, where $R^2$ is the variance explained by prespecified pre-assignment covariates, the standardised variance $v$ is $2(1-R^2)/n$ for the two-arm contrast and $(1-R^2)/n$ for the matching contrast, with $n$ retained per cell. The known-variance, two-sided normal-theory power is
\[
\pi(d,v)=\Phi(|d|/\sqrt{v}-z_{1-\alpha/2})
+\Phi(-|d|/\sqrt{v}-z_{1-\alpha/2}).
\]
Here $\Phi$ is the standard normal cumulative distribution function and $z_p$ its $p$ quantile. Recruitment of $K\lceil n/(1-\ell)\rceil$ participants across $K$ cells inflates the target for assumed loss $\ell$; it does not guarantee retention or correct informative missingness.

\begin{table}[htbp]
\centering\small
\caption{Illustrative total recruitment for a single primary contrast: 80\% power, two-sided $\alpha=.05$, 15\% assumed loss.}
\label{tab:power}
\begin{tabular}{@{}rrrrr@{}}
\toprule
&\multicolumn{2}{c}{Two arms}&\multicolumn{2}{c}{Four cells}\\
Effect & $R^2=0$ & $R^2=.25$ & $R^2=0$ & $R^2=.25$\\
\midrule
.20 & 926 & 696 & 928 & 700\\
.30 & 412 & 310 & 416 & 312\\
.40 & 234 & 176 & 236 & 176\\
.50 & 150 & 114 & 152 & 116\\
\bottomrule
\end{tabular}
\end{table}

The companion Python script computes the smallest retained cell size meeting the target and checks the formula by simulating independent Gaussian cell means. At a standardised effect of .30 and $R^2=0$, 175 retained participants per arm and 88 per factorial cell yield simulated powers of .80065 and .80200, respectively, in 20,000 repetitions per design. At the same sample sizes under the null, rejection proportions are .04865 and .04925 (20,000 repetitions each). Monte Carlo standard errors for the power estimates are approximately .0028.

The sensitivity analysis assumes known variance. For a balanced Gaussian design with common estimated variance, the corresponding noncentral-$t$ distribution has $K(n-1)$ degrees of freedom; at $d=.30$ and $R^2=0$, the minimum is 176 per arm or 88 per factorial cell. A supplementary Python diagnostic uses Welch--Satterthwaite intervals with estimated cell variances, unequal allocation, heterogeneous variance and independent retention; it reports bias, rejection, coverage, valid replicate counts and Monte Carlo uncertainty. Separate synthetic scenarios examine clustered diary slopes, expertise-by-interface interactions and enjoyment trends. Extended mixed-likelihood diagnostics (2,000 replicates per scenario; 80 participants, eight observations) yielded H6 null rejection .052 without missingness and .0815 under MNAR. H10 power at an illustrative contrast of $-.35$ was .193 without missingness and .133 under MAR; covariance boundary solutions were frequent. These findings preclude confirmatory H10 at that size. Revised synthetic diagnostics (2,000 replicates per cell) support conservative H6 CR2 inference under MAR (null rejection .037 at $\alpha=.04$), not unrestricted MNAR (.067). With active final-session recovery (approximately 83\% observed under MAR), the revised endpoint estimator at 800 assigned participants yields .868 power for an illustrative contrast of $-.35$; 600 yields .730. This is a conditional planning scenario, not an empirically justified recruitment target. Recalculate power and calibrate inference for the selected estimator, clustering, attrition and multiplicity. Choose the smallest effect of interest from outcome meaning, not pooled effects of other musical activities. Begin with qualitative discovery, cognitive interviews and feasibility work before a focused confirmatory experiment. Pilot data should inform measurement properties, variance and completion rates, not determine the target effect through an unstable preliminary estimate.


\paragraph{Rounded recruitment sensitivity (27 November 2026).}
For the illustrative factorial matching contrast with $d_\psi=.30$, no baseline variance reduction and 15\% independent loss, the proposed $N=420$ allocates 105 participants per cell and yields 357 expected observed outcomes. The noncentral-$t$ approximation gives power .80689. An independent local reproduction generated 100,000 Gaussian cell-summary datasets at each of the null and alternative, with binomial retention, estimated cell variances and Welch--Satterthwaite intervals. Null rejection was .05084; alternative rejection was .80578 (Monte Carlo SE .00125); coverage was .94916 and .94982, respectively. Mean retained $N$ was 356.93 and 357.02. These calculations do not calibrate bounded outcomes, informative loss or the remaining hypotheses. The recruitment target remains illustrative.


\subsection{Joint prediction calibration}
A separate synthetic screen evaluated the intersection-union decision for $s_0>0$ and $s_1>0$. Gaussian independent cell summaries were generated with binomial outcome-independent retention and estimated variances. The grid crosses six operational sample sizes (420, 600, 800, 1000, 1200, 1400), eight symmetric, asymmetric, global-null, partial-null and reversed simple-effect configurations, and two variance/retention profiles (20,000 independent datasets per cell, 1,920,000 total; seed 20261203). The equal profile uses unit residual standard deviations and .85 retention; the heterogeneous profile uses standard deviations (.8, 1.2, 1.2, .8) and retention (.90, .80, .75, .90) in cell order 00, 01, 10, 11. All simulated fits were valid. At $s_0=s_1=.30$ in common unit-scale points, joint rejection probabilities for $N=420,800,1000,1200$ were .40685, .75240, .85760 and .91910 in the equal profile, and .37655, .70515, .81770 and .88940 in the heterogeneous profile. For $N=420$, an asymmetric (.15,.45) configuration had joint power .23660 despite mean-contrast rejection .80730. A partial-null (0,.60) configuration produced joint rejection .04970, whereas mean-contrast rejection was .80910. Across the simulated union-null configurations the largest joint rejection rate was .05225, within Monte Carlo uncertainty of .05; simple-effect pointwise coverage ranged .94525 to .95385 across all cells. At $N=1000$, joint power Monte Carlo standard errors were .00247 and .00273 for the two symmetric profiles. These simulations do not establish that the bounded rating is Gaussian, invariant or missing at random. The symmetric-effect grid identifies $N=1000$ as the first evaluated candidate exceeding .80 in both profiles; it does not establish a minimum for asymmetric effects or a recruitment target. Figure~\ref{fig:cp-power} displays the reported subset of joint-power estimates and contrasts the symmetric, asymmetric and partial-null configurations at $N=420$. The mean-contrast rejection probability is not the power of CP. Preserve the earlier mean-contrast results as a distinct historical estimand. Pilot evidence, justified component effect sizes, measurement comparability, informative-missingness analysis and owner approval remain necessary. The local execution used Python numerical dependencies NumPy 2.3.5 and SciPy 1.17.0; repository-pinned CI reproduction remains separate.

\begin{figure}[htbp]
\centering
\begin{tikzpicture}[x=1cm,y=1cm,font=\small]
\begin{scope}[xscale=0.006,yscale=4]
\draw[->] (380,0)--(1245,0) node[below right] {$N$};
\draw[->] (380,0)--(380,1.05);
\foreach \v in {0,.2,.4,.6,.8,1}{\draw[gray!25] (380,\v)--(1220,\v);\node[left] at (380,\v) {\v};}
\foreach \n in {420,800,1000,1200}{\node[below] at (\n,0) {\n};}
\draw[dashed] (380,.8)--(1220,.8);
\draw[thick,blue!70!black] plot coordinates {(420,.40685)(800,.75240)(1000,.85760)(1200,.91910)};
\draw[thick,orange!85!black,densely dotted] plot coordinates {(420,.37655)(800,.70515)(1000,.81770)(1200,.88940)};
\foreach \n/\v in {420/.40685,800/.75240,1000/.85760,1200/.91910}{\fill[blue!70!black] (\n,\v) circle[radius=1.5pt];}
\foreach \n/\v in {420/.37655,800/.70515,1000/.81770,1200/.88940}{\draw[orange!85!black,fill=white] (\n,\v) circle[radius=1.5pt];}
\node[above,align=center,font=\small\bfseries] at (800,1.05) {A. Joint power};
\node[above right,blue!70!black,font=\scriptsize] at (420,.45) {Equal};
\node[below right,orange!85!black,font=\scriptsize] at (420,.32) {Heterogeneous};
\end{scope}
\begin{scope}[xshift=8.1cm,x=1cm,y=4cm]
\draw[->] (0,0)--(4.7,0);
\draw[->] (0,0)--(0,1.05);
\foreach \v in {0,.2,.4,.6,.8,1}{\draw[gray!25] (0,\v)--(4.5,\v);\node[left] at (0,\v) {\v};}
\fill[blue!70!black] (.65,0) rectangle (1.25,.23660);
\fill[gray!55] (1.28,0) rectangle (1.88,.80730);
\fill[blue!70!black] (2.65,0) rectangle (3.25,.04970);
\fill[gray!55] (3.28,0) rectangle (3.88,.80910);
\node[below,align=center,font=\scriptsize] at (1.25,0) {(.15,.45)};
\node[below,align=center,font=\scriptsize] at (3.25,0) {(0,.60)};
\node[above,font=\small\bfseries] at (2.3,1.05) {B. Rejection at $N=420$};
\node[blue!70!black,font=\scriptsize] at (1.05,.95) {Joint};
\node[font=\scriptsize] at (3.3,.95) {Mean contrast};
\end{scope}
\end{tikzpicture}

\caption{Reported synthetic calibration results. Panel A: joint rejection probabilities for symmetric simple effects (.30, .30) at four reported sample sizes. Lines connect evaluated points only; the dashed line is the .80 planning threshold. Panel B: joint and mean-contrast rejection probabilities at $N=420$ for asymmetric and partial-null configurations. Under the partial null the joint rejection probability is a Type I error rate, not power. No participant observations are represented.}
\label{fig:cp-power}
\end{figure}

\subsection{Bounded-outcome and observation sensitivity}
A supplementary synthetic design diagnostic represents the candidate rating as an integer from 0 to 10 using a beta-binomial distribution. The full-population simple effects are defined on the observed score scale, not a latent Gaussian scale. Low, intermediate and high baseline means assess floor and ceiling behaviour; two concentration parameters vary score dispersion. A binary pre-assignment prognostic factor influences both the score and observation under MAR, whereas MNAR additionally depends on the unobserved score. The diagnostic compares available-case Welch inference with equally weighted stratum means and Satterthwaite uncertainty under a known target stratum distribution. Standardisation addresses the specified MAR mechanism under adequate stratum support, but does not identify MNAR outcomes. The planned grid crosses three sample sizes, five effect patterns, three baseline means, four observation mechanisms and two dispersion settings. Report unconditional conjunction rejection, valid and failed counts, Monte Carlo standard errors and conditional bias, RMSE and coverage. Results are not reported here until exact-source execution is verified; the diagnostic does not establish a validated rating, a recruitment target or an approved confirmatory estimator.

\printbibliography
\end{document}
